\section*{الفصل \ref{ch:b3:quantitative-geomorphology} --- الجيومورفولوجيا الكمّية}

\begin{solution}{exo:b3:quantitative-geomorphology:1}
$a/\sqrt{\pi\kappa t} = 0.30 - 0.04 = 0.26$، إذن $\kappa t = (2/0.26)^2/\pi =
\qty{18.8}{m^2}$، و$t = 18.8/(5\times10^{-3}) = \qty{3.8}{kyr}$، ضمن
عامل قدره اثنان في الاتجاهين، نظرًا لمدى $\kappa$.
\end{solution}

\begin{solution}{exo:b3:quantitative-geomorphology:2}
$E = \kappa\,|\nabla^2 z| = 4\times10^{-3}\times0.012 =
\qty{4.8e-5}{m/yr}$، أي \qty{0.048}{mm/yr}. التضاريس $E L^2/(2\kappa) =
0.012\times80^2/2 = \qty{38.4}{m}$.
\end{solution}

\begin{solution}{exo:b3:quantitative-geomorphology:3}
$E = 2\times10^{-6}\times(25\times10^6)^{0.5}\times0.01 = 2\times10^{-6}
\times5000\times0.01 = \qty{1.0e-4}{m/yr}$، أي \qty{0.1}{mm/yr}. بما أنّ $E$
بالوحدة \unit{m/yr}، و$A^m$ بالوحدة $\unit{m}^{2m}$، و$S$ بلا وحدة، فإنّ $K$ بالوحدة
$\unit{m}^{1-2m}\,\unit{yr^{-1}}$: \unit{yr^{-1}} عندما $m = 0.5$.
\end{solution}

\begin{solution}{exo:b3:quantitative-geomorphology:4}
(أ) ثابت: يزيد الصخر الصلب انحدار النهر حيث يظهر، ويبقى الانحدار المفاجئ
مرتبطًا بالصخر الأساسي. (ب) مثبّت عند الصدع: ما دام الصدع ينزلق،
يزداد انحدار النهر عند عبوره له. (ج) متحرّك: أدّى رفع
الساحل إلى خفض \omterm{def:b3:quantitative-geomorphology:knickpoint}{مستوى القاعدة} لكل الأنهار دفعةً واحدة؛ وتنتقل نقاط الانقطاع
نحو المنبع، وكما تتنبّأ مبرهنة الموجة الحركية عندما $n
= 1$، فإنّها تقع على كنتور واحد.
\end{solution}

\begin{solution}{exo:b3:quantitative-geomorphology:5}
$\theta = m/n = 0.5$. $S = (U/K)^{1/2}A^{-1/2} = (10^{-4}/10^{-9})^{1/2}
\times(10^7)^{-1/2} = 316\times3.16\times10^{-4} = 0.10$. مضاعفة $U$
تضرب $k_s = (U/K)^{1/n}$ في $\sqrt2 = 1.41$؛ وعندما $n = 1$ تضربه في 2:
وكلّما زادت $n$، قلّت استجابة الانحدار للرفع.
\end{solution}

\begin{solution}{exo:b3:quantitative-geomorphology:6}
بقياس $x$ من المخرج نحو المنبع، $\mathrm{d}x = -\mathrm{d}d$، ومنه $\chi(d)
= \int_d^L (A_0/c\,d'^{h})^{\theta}\,\mathrm{d}d' = A_0^{\theta} c^{-\theta}
\int_d^L d'^{-h\theta}\,\mathrm{d}d' = A_0^{\theta}c^{-\theta}(L^p -
d^p)/p$. هنا $p = 0.19$، و$A_0^{\theta}c^{-\theta} = 10^{2.7}/2^{0.45}$،
و$\chi = \qty{2935}{m}$.
\end{solution}

\begin{solution}{exo:b3:quantitative-geomorphology:7}
$c = 2\times10^{-6}\times(10^8)^{0.5} = \qty{0.02}{m/yr}$, \qty{20}{m/kyr};
تستغرق المسافة \qty{5}{km} بتلك السرعة \qty{250}{kyr}. وتقلّ المساحة، ومعها
سرعةُ الموجة تتناقص نحو المنبع، فتتباطأ \omterm{def:b3:quantitative-geomorphology:knickpoint}{نقطة الانقطاع} كلّما صعدت.
\end{solution}

\begin{solution}{exo:b3:quantitative-geomorphology:8}
$\bar e = 0.40\times600 = \qty{240}{m}$؛ ترتفع القمم بمقدار $0.80\times240 =
\qty{192}{m}$؛ وينخفض متوسط الارتفاع بمقدار $240 - 192 = \qty{48}{m}$.
\end{solution}

\begin{solution}{exo:b3:quantitative-geomorphology:9}
$h = h_0\ln(P_0/E)$: $0.5\ln 6 = \qty{0.90}{m}$ و$0.5\ln1.5 =
\qty{0.20}{m}$. عند \qty{0.4}{mm/yr} تتجاوز التعرية أعلى معدّل للإنتاج،
$P_0$: تُجرَّد التربة ويصبح المنحدر صخرًا عاريًا، ويكون خفضه
محدودًا بالتجوية.
\end{solution}

\begin{solution}{exo:b3:quantitative-geomorphology:10}
في \omterm{def:b3:quantitative-geomorphology:steady-profile}{الحالة المستقرّة} $\mathrm{d}q/\mathrm{d}x = E$ و$q = 0$ عند خط التقسيم، ومنه
$q = E x$. ثمّ $\kappa S = E x\,(1 - S^2/S_c^2)$، وهي معادلة تربيعية
جذرها الموجب $S = [-\kappa + \sqrt{\kappa^2 + 4E^2x^2/S_c^2}\,]/(2Ex/S_c^2)$. عند
20، 50، 100 و\qty{200}{m}: 0.36، 0.68، 0.89 و1.03، مقابل 0.4، 1.0،
2.0 و4.0 للانتشار الخطّي ($S = Ex/\kappa$). للتلّ قمّةٌ
مستديرة، ثمّ جوانب يميل تدرّجها نحو $S_c$: منحدرات شبه مستوية، ذات
تضاريس أقلّ بكثير من تضاريس القطع المكافئ.
\end{solution}

\begin{solution}{exo:b3:quantitative-geomorphology:11}
بمبرهنة الموجة الحركية يبلغ التغيّر نقطةً ذات الإحداثي
$\chi$ بعد $\chi/(KA_0^m)$؛ وتكون الشبكة قد استجابت استجابةً تامّة حين تبلغ
آخر رأس للقناة، ذا أكبر $\chi$. $K A_0^m = 10^{-6}\times
1000 = \qty{e-3}{m/yr}$: يستجيب الجذع خلال \qty{5.46}{Myr}، والشبكة
خلال \qty{7.24}{Myr}، من رتبة مليون سنة متوقّعة.
\end{solution}

\begin{solution}{exo:b3:quantitative-geomorphology:12}
فوق الملتقى يكون التدرّج المستقرّ للرافد $S = (U/2K)A^{-\theta}$،
ومنه $\mathrm{d}z/\mathrm{d}\chi = (U/2K)A_0^{-\theta}$، أي نصف تدرّج الجذع.
تبلغ \omterm{def:b3:quantitative-geomorphology:knickpoint}{نقطة الانقطاع} الملتقى حين يكون الجذع عند $\chi_j$؛ وبعد
زمنٍ إضافي $\Delta t$ ينتقل الجذع إلى $\chi_k$ حيث $\chi_k -
\chi_j = K A_0^m\,\Delta t$، بينما ينتقل الرافد بسرعة $2KA^m$،
إلى $\chi$ حيث $\chi - \chi_j = 2KA_0^m\,\Delta t = 2(\chi_k -
\chi_j)$.
\end{solution}

\begin{solution}{pb:b3:quantitative-geomorphology:1}
\textbf{1.} $4.5\times12000^{1.7} = \qty{3.87e7}{m^2}$, \qty{38.7}{km^2};
$4.5\times300^{1.7} = \qty{7.3e4}{m^2}$, \qty{0.073}{km^2}.

\textbf{2.} $\theta \approx -\log(0.007855/0.06641)/\log(37.08/0.388) =
0.928/1.980 = 0.47$.

\textbf{3.} نعم: مع تشتّت نحو \qty{10}{\percent} على اثنتي عشرة نقطة،
$0.466$ لا يختلف اختلافًا معنويًا عن 0.5.

\textbf{4.} $0.06641\times\sqrt{388\,000} = \qty{41.4}{m}$;
$0.007855\times\sqrt{3.708\times10^7} = \qty{47.8}{m}$. يقدّر المتوسّط الهندسي،
\qty{45.1}{m}، يقدّر $k_s$ من أجل $\theta = 0.5$.

\textbf{5.} التدرّج هو الفرق بين ارتفاعين، يحمل كلٌّ منهما عدم يقين
بمقدار ديسيمترات إلى أمتار، مقسومًا على مسافة قصيرة؛ وعلى المقاطع اللطيفة
يمثّل الخطأ جزءًا كبيرًا من الفرق. أمّا $\chi$ فيُكامِل بدل
أن يُفاضِل.

\textbf{6.} $p = 1 - 1.7\times0.5 = 0.15$; $\chi = (1000/\sqrt{4.5})\,
(12000^{0.15} - 9000^{0.15})/0.15 = \qty{543.1}{m}$.

\textbf{7.} $(193.4 - 170)/543.1 = 0.0431$; $k_{sn} = 0.0431\times A_0^{0.5}
= \qty{43.1}{m}$.

\textbf{8.} تقدير كاي: يعتمد على الارتفاعات فوق \qty{3}{km} من
النهر، لا على اثني عشر تدرّجًا مشوّشًا.

\textbf{9.} $K = U/k_{sn} = 4.3\times10^{-5}/43.1 = \qty{1.00e-6}{yr^{-1}}$.

\textbf{10.} $10^{-6}\times\sqrt{3.87\times10^7}\times0.00691 =
\qty{4.30e-5}{m/yr}$، \qty{43}{m/Myr}: يواكب المخرج
الرفع.

\textbf{11.} عند المخرج $10^{-6}\times6222 = \qty{6.2}{m/kyr}$؛
وعند الشلّال، $d = \qty{7.40}{km}$، $A = \qty{1.70e7}{m^2}$
و\qty{4.1}{m/kyr}.

\textbf{12.} $170 + 43\times(4.5)^{-0.5}(12000^{0.15} - 300^{0.15})/0.15 =
\qty{405}{m}$: التضاريس بمقدار \qty{235}{m}.

\textbf{13.} من أجل $n = 1$ تكون المعادلة موجةً حركيةً: نقاط على جانبي
الدرجة ترتفع بمقدار $U$ في وحدة الزمن، فيبقى ارتفاعها،
أي مقدار السقوط \qty{25}{m}.

\textbf{14.} $\chi = (1000/\sqrt{4.5})(12000^{0.15} - 7400^{0.15})/0.15 =
\qty{899.4}{m}$.

\textbf{15.} $t = \chi/(KA_0^{0.5}) = 899.4/10^{-3} = \qty{899}{kyr}$.

\textbf{16.} عند $\chi = 899.4 + 300 = \qty{1199}{m}$:
$d^{0.15} = 12000^{0.15} - 1199\times0.15\times\sqrt{4.5}/1000$, $d =
\qty{6.25}{km}$، \qty{5.75}{km} أعلى المجرى من المخرج.

\textbf{17.} عند الملتقى ($d = \qty{9.5}{km}$) $\chi = \qty{442.8}{m}$؛
يقع الشلّال حيث يبلغ $\chi$ الخاص بالرافد \qty{899.4}{m}:
$s^{0.15} = 5000^{0.15} - (899.4 - 442.8)\times0.15\times\sqrt{1.5}/1000$,
$s = \qty{4.27}{km}$، \qty{0.73}{km} فوق الملتقى.

\textbf{18.} تقع ملتقياتها عند $\chi = 1954$ و\qty{999}{m}، وهي تتجاوز
\qty{899}{m}: لم تبلغها \omterm{def:b3:quantitative-geomorphology:knickpoint}{نقطة الانقطاع} بعدُ، فالجذع هناك
ما زال على مستواه القديم.

\textbf{19.} $E = 5\times10^{-3}\times0.0086 = \qty{4.3e-5}{m/yr}$,
\qty{43}{m/Myr}: تتآكل التلال فوق الشلّال بمعدّل الرفع، في
\omterm{def:b3:quantitative-geomorphology:steady-profile}{الحالة المستقرّة}.

\textbf{20.} $150^2/(5\times10^{-3}) = \qty{4.5}{Myr}$، خمسة أضعاف عمر
الشلّال.

\textbf{21.} خُفِضت قاعدتها بمقدار \qty{25}{m} مع القناة: فأجزاؤها
السفلى شديدة الانحدار، وربما تنزلق، بينما تستمرّ قممها في
التآكل بالمعدّل القديم. وتصعد موجةُ ازدياد الانحدار كلّ سفح؛ وبعد بضعة
ملايين من السنين تعود إلى القطع المكافئ نفسه، أدنى بمقدار \qty{25}{m}.

\textbf{22.} يتغيّر $t = \chi k_{sn}/(U A_0^{0.5})$ تبعًا، أي يتناسب مع $1/U$: من
$899/1.2 = 749$ إلى $899/0.8 = \qty{1124}{kyr}$.

\textbf{23.} مع $n = 2$ تنمو سرعة الموجة مع التدرّج: فالوجه الشديد
يتحرّك أسرع من المقاطع المحيطة به فيبقى حادًّا، لكنّ
سرعته والعمر المستنتَج منه سيختلفان.

\textbf{24.} انخفض \omterm{def:b3:quantitative-geomorphology:knickpoint}{مستوى القاعدة} لنهر فاردن
قبل \textbf{\qty[text-series-to-math=true]{900}{kyr}} (\qty{899}{kyr} وفق
الحساب).
\end{solution}
